\documentclass[11pt]{article}
\usepackage{amsmath,amssymb,amsthm}
\usepackage[margin=1.2in]{geometry}
\usepackage{booktabs}
\usepackage{lmodern}
\usepackage[T1]{fontenc}
\usepackage{listings}
\usepackage[htt]{hyphenat}
\usepackage{hyperref}
\emergencystretch=3em
\sloppy

\lstset{basicstyle=\ttfamily\footnotesize,breaklines=true,columns=fullflexible,
        xleftmargin=1.5em,frame=none,extendedchars=true,inputencoding=utf8,
        literate={σ}{{$\sigma$}}1 {λ}{{$\lambda$}}1
        {Δ}{{$\Delta$}}1 {ℝ}{{$\mathbb{R}$}}1 {ℂ}{{$\mathbb{C}$}}1 {ℕ}{{$\mathbb{N}$}}1
        {→}{{$\to$}}1 {↦}{{$\mapsto$}}1 {∀}{{$\forall$}}1 {∃}{{$\exists$}}1
        {≤}{{$\le$}}1 {≥}{{$\ge$}}1 {≠}{{$\ne$}}1 {∈}{{$\in$}}1 {∫}{{$\int$}}1
        {α}{{$\alpha$}}1 {ι}{{$\iota$}}1 {κ}{{$\kappa$}}1 {ρ}{{$\rho$}}1 {ξ}{{$\xi$}}1
        {Ξ}{{$\Xi$}}1 {Ψ}{{$\Psi$}}1 {·}{{$\cdot$}}1 {ᴴ}{{$^{H}$}}1
        {₀}{{$_0$}}1 {₁}{{$_1$}}1 {₂}{{$_2$}}1 {∧}{{$\wedge$}}1 {∨}{{$\vee$}}1
        {¬}{{$\neg$}}1 {↔}{{$\leftrightarrow$}}1 {ℤ}{{$\mathbb{Z}$}}1
        {∑}{{$\sum$}}1 {∏}{{$\prod$}}1 {√}{{$\surd$}}1 {⁺}{{$^{+}$}}1
        {≡}{{$\equiv$}}1 {∘}{{$\circ$}}1 {⟨}{{$\langle$}}1 {⟩}{{$\rangle$}}1}

\newtheorem{theorem}{Theorem}
\newtheorem{proposition}[theorem]{Proposition}
\newtheorem{lemma}[theorem]{Lemma}
\newtheorem{corollary}[theorem]{Corollary}
\theoremstyle{definition}
\newtheorem{definition}[theorem]{Definition}
\newtheorem{remark}[theorem]{Remark}

\DeclareMathOperator{\Imop}{Im}
\DeclareMathOperator{\Reop}{Re}

\title{De Bruijn's strip-narrowing theorem for the Riemann \texorpdfstring{$\Xi$}{Xi}-function,\\
       formalized in Lean 4}
\author{Samuel Lavery}
\date{Draft --- \today}

\begin{document}
\maketitle

\begin{abstract}
We report a complete, machine-checked proof in Lean 4 (Mathlib) of de Bruijn's 1950 theorem on the
heat-flow deformations of the Riemann $\Xi$-function: writing
$\Xi_\lambda(z)=\int_{\mathbb R} e^{\lambda u^{2}}K(u)e^{izu}\,du$ with $\Xi_0=\Xi$, every $\Xi_\lambda$
with $\lambda\ge\tfrac18$ has only real zeros.  In the normalization standard in the
de Bruijn--Newman literature this is the bound $\Lambda\le\tfrac12$ on the de Bruijn--Newman constant.
The development is self-contained --- Mathlib contains no Laguerre--P\'olya theory, no preservation of
real-rootedness under locally uniform limits, no Hadamard product for $\xi$, and no de Bruijn--Newman
material --- so the Hadamard factorization of $\xi$, the polynomial strip-narrowing step, the passage
from imaginary shifts to the heat semigroup, and a Hurwitz limit theorem are all built here.  The
result is $216$ theorems over $14$ modules and $3462$ lines, with axiom footprint
$\{\texttt{propext},\ \texttt{Classical.choice},\ \texttt{Quot.sound}\}$ and no \texttt{sorry}.
We also record a transversality theorem at the polynomial level --- every real zero of the shift of a
real-rooted polynomial is simple --- together with an explicit account of why it does \emph{not}
descend to $\Xi_\lambda$.
The Riemann Hypothesis is not proved and is not approached: the interval $\lambda\in(0,\tfrac18)$
remains, and closing it is equivalent to $\Lambda\le 0$, which is RH.  We close by reporting a
numerical experiment locating where the $\tfrac18$ is slack.
\end{abstract}

\section{What is formalized}

Let $\xi$ be the completed Riemann zeta function and put
\[
  \Xi(z)\;=\;\xi\!\left(\tfrac12+iz\right).
\]
Then $\Xi$ is entire, real on $\mathbb R$, and its zeros are the images of the nontrivial zeros of
$\zeta$: if $\rho=\beta+i\gamma$ is a nontrivial zero, the corresponding $z$ is
$z=\gamma-i(\beta-\tfrac12)$, so
\begin{equation}\label{eq:strip}
  \Xi(z)=0 \;\Longrightarrow\; |\Imop z|=\left|\beta-\tfrac12\right|<\tfrac12 ,
\end{equation}
because $0<\beta<1$.  The Riemann Hypothesis is the assertion that every zero of $\Xi$ is real.

$\Xi$ has a classical Fourier representation $\Xi(z)=\int_{\mathbb R}K(u)e^{izu}\,du$ \cite[\S10.1]{Titchmarsh} with an explicit even,
positive, rapidly decaying kernel $K$ built from the Jacobi theta function (\S\ref{sec:normal}).
Deforming the kernel by a Gaussian gives the \emph{flow}
\begin{equation}\label{eq:flow}
  \Xi_\lambda(z)\;=\;\int_{\mathbb R} e^{\lambda u^{2}}K(u)\,e^{izu}\,du ,
  \qquad \Xi_0=\Xi .
\end{equation}
This is the backward heat flow: since $\partial_z^{2}e^{izu}=-u^{2}e^{izu}$, we have
$\Xi_\lambda=e^{-\lambda\partial_z^{2}}\Xi$ formally, and the zeros of $\Xi_\lambda$ migrate toward the
real axis as $\lambda$ increases.  The object of this paper is the following theorem, proved by
de Bruijn in 1950 (\cite{deBruijn1950}, Thm.~13) and here formalized.  De Bruijn's statement is the
strip contraction in the form $|\Imop z|\le\bigl(\max(\Delta^{2}-\lambda,0)\bigr)^{1/2}$ for the kernel
factor $e^{\lambda t^{2}/2}$, which with the factor $e^{\lambda u^{2}}$ used here reads
$\bigl(\max(\Delta^{2}-2\lambda,0)\bigr)^{1/2}$ --- the law appearing as
Proposition~\ref{prop:shift} below.

\begin{theorem}[\texttt{flowFrontier\_one\_eighth}]\label{thm:main}
For every $\lambda$ with $\tfrac18\le\lambda\le 1$, every zero of $\Xi_\lambda$ is real.
\end{theorem}

In Lean, with \texttt{flowStage} the function \eqref{eq:flow}:

\begin{lstlisting}
def FlowFrontier (lam₀ : ℝ) : Prop :=
  ∀ lam : ℝ, lam₀ ≤ lam → lam ≤ 1 → ∀ z : ℂ, flowStage lam z = 0 → z.im = 0

theorem flowFrontier_one_eighth : FlowFrontier (1 / 8)
\end{lstlisting}

\noindent The upper cut $\lambda\le1$ is a convenience of the ambient development, not a restriction of
the method; the narrowing argument is monotone in $\lambda$ and \texttt{FlowFrontier.mono} transports
the conclusion to every $\lambda_0\ge\tfrac18$.

The mechanism has three parts, and each is elementary in isolation:

\begin{enumerate}
\item[(i)] a \emph{polynomial} step: the map $p\mapsto p(\cdot+ih)+p(\cdot-ih)$ carries a real
  polynomial whose zeros lie in $(\Imop w)^{2}\le\sigma$ to one whose zeros lie in
  $(\Imop w)^{2}\le\max(\sigma-h^{2},0)$ (\S\ref{sec:core}--\S\ref{sec:class});
\item[(ii)] an \emph{approximation} step: $\Xi$ is a locally uniform limit of real polynomials with all
  zeros in $(\Imop w)^{2}\le\tfrac14$, supplied by the Hadamard product (\S\ref{sec:hadamard});
\item[(iii)] a \emph{limit} step: iterating (i) with $N$ shifts of size $h$ where $Nh^{2}=2\lambda$
  realizes the Gaussian $e^{\lambda u^{2}}$ in the limit, and real-rootedness survives locally uniform
  limits by Hurwitz (\S\ref{sec:cosh}--\S\ref{sec:assembly}).
\end{enumerate}

\section{Normalization}\label{sec:normal}

The de Bruijn--Newman literature is not uniform in its scaling, and the numerical value of the
threshold depends on it, so we fix conventions explicitly.

Let $\psi(x)=\sum_{n\ge1}e^{-\pi n^{2}x}$ and set $\Psi(u)=\psi(e^{2u})$.  The functional equation
$\psi(1/x)=\sqrt{x}\,\psi(x)+\tfrac{\sqrt x-1}{2}$ becomes
\begin{equation}\label{eq:reflect}
  \Psi(u)=e^{-u}\Psi(-u)+\tfrac{e^{-u}-1}{2},
\end{equation}
which is compiled as \texttt{Psi\_reflect} (from Mathlib's \texttt{evenKernel\_functional\_equation}
at $a=0$; in the formalization $\Psi$ is the restriction to $\mathbb R$ of
$\texttt{PsiC}(w)=\tfrac12(\vartheta_2(0,ie^{2w})-1)$, holomorphic on $|\Imop w|<\pi/4$).  The kernel is
\[
  K(u)\;=\;e^{u/2}\bigl(\Psi''(u)+\Psi'(u)\bigr),
\]
and differentiating \eqref{eq:reflect} twice gives $\Psi''(u)+\Psi'(u)=e^{-u}(\Psi''(-u)+\Psi'(-u))$,
hence $K(-u)=K(u)$: the kernel is even (\texttt{Kker\_even}).  That $\int_{\mathbb R}K(u)e^{izu}du=\Xi(z)$
is \texttt{xiSection\_eq\_Ktrans}, proved by analytic continuation from the half-plane
$\Imop z<-\tfrac12$ where the integral converges absolutely term by term.

With this $u$, the strip bound \eqref{eq:strip} reads $\Delta=\tfrac12$, and Theorem~\ref{thm:main}
gives real-rootedness at $\lambda\ge\tfrac18$.  The de Bruijn--Newman constant $\Lambda$ is defined in
the literature through the family $H_t$ normalized so that $H_0(z)=\tfrac18\Xi(z/2)$; halving the
$z$-variable doubles the half-width to $\Delta_{\text{std}}=1$ and divides the Gaussian exponent by
four, so
\begin{equation}\label{eq:norm}
  \lambda \;=\; t/4 .
\end{equation}
Relation \eqref{eq:norm} is not ours: Ki--Kim--Lee state it explicitly, defining the de Bruijn--Newman
constant as $\Lambda=4\lambda(0)$ with $\lambda(0)$ the infimum of $\lambda$ for which $\Xi_\lambda$ has
only real zeros \cite{KKL2009}.  Theorem~\ref{thm:main} is therefore exactly de Bruijn's published
bound $\Lambda\le\tfrac12$ \cite{deBruijn1950}.  Two further consistency checks: de Bruijn's
$t=\tfrac12$ maps to $\lambda=\tfrac18$, and the Polymath~15 bound $t\le0.22$
\cite{Polymath15} maps to $\lambda\le0.055$, correctly strictly better than $\tfrac18$.  A reader who
prefers the $H_t$ normalization should read every $\lambda$ below as $t/4$.

\section{The elementary core: one identity}\label{sec:core}

Everything in part (i) rests on a single algebraic identity, which we state first because it is the
whole mechanism and it is two lines.

For real $a,b$ let $Q_{a,b}(z)=(z-a)^{2}+b^{2}$, the real quadratic with roots $a\pm ib$.

\begin{lemma}[\texttt{pairFactor\_normSq\_shift}]\label{lem:core}
For all real $a,b,h$ and all $z\in\mathbb C$,
\[
  \bigl|Q_{a,b}(z+ih)\bigr|^{2}-\bigl|Q_{a,b}(z-ih)\bigr|^{2}
  \;=\;8\,(\Imop z)\,h\,\Bigl((\Reop z-a)^{2}+(\Imop z)^{2}+h^{2}-b^{2}\Bigr).
\]
\end{lemma}

\begin{proof}
Expand both sides in $\Reop z$, $\Imop z$ and compare; in Lean this is \texttt{simp} on the real and
imaginary parts followed by \texttt{ring}.
\end{proof}

The consequence is immediate and is the engine of the whole paper.

\begin{corollary}[\texttt{pairFactor\_norm\_shift\_lt}]\label{cor:core}
If $\Imop z>0$, $h>0$ and $b^{2}<(\Imop z)^{2}+h^{2}$, then
$\bigl|Q_{a,b}(z-ih)\bigr|<\bigl|Q_{a,b}(z+ih)\bigr|$.
\end{corollary}

\noindent In words: \emph{a conjugate pair of roots that is not too high up sees the upward shift as
strictly larger than the downward one}, uniformly in the real part.  A real polynomial $p$ factors
over $\mathbb C$ into such quadratics together with real linear factors, and a real linear factor
$z-a$ is the case $b=0$ of $Q_{a,b}$ after squaring; taking the product of strict inequalities over
all roots (\texttt{polynomial\_norm\_shift\_lt}) gives

\begin{proposition}[\texttt{polynomial\_shift\_zero\_strip}]\label{prop:polystep}
Let $p\in\mathbb R[X]$, $p\ne0$, $h>0$, and suppose every root $w$ of $p$ satisfies
$(\Imop w)^{2}\le\Delta^{2}$.  If $p(z+ih)+p(z-ih)=0$ then
\[
  \Imop z=0 \qquad\text{or}\qquad (\Imop z)^{2}+h^{2}\le\Delta^{2}.
\]
\end{proposition}

\begin{proof}
If $\Imop z>0$ and $(\Imop z)^{2}+h^{2}>\Delta^{2}$ then Corollary~\ref{cor:core} applies to every
root, so $\|p(z-ih)\|<\|p(z+ih)\|$; but $p(z+ih)=-p(z-ih)$ forces equality of norms, a contradiction.
The case $\Imop z<0$ follows by conjugating, using that $p$ has real coefficients
($\texttt{Polynomial.aeval\_conj}$), which exchanges the two shifted arguments.  No simplicity
assumption on the roots is used: multiplicities are carried by the multiset of roots throughout.
\end{proof}

We emphasise what Proposition~\ref{prop:polystep} costs: a single \texttt{ring} identity, a product of
strict inequalities, and the triangle inequality.  There is no potential theory and no
Laguerre--P\'olya theory in it.

\section{The strip class and the shift}\label{sec:class}

\begin{definition}[\texttt{StripClass}]
For $\sigma\in\mathbb R$ and $f:\mathbb C\to\mathbb C$, say $f$ lies in the \emph{strip class at
$\sigma$} if there are nonzero real polynomials $P_n$ with every root of every $P_n$ satisfying
$(\Imop w)^{2}\le\sigma$, such that $P_n\to f$ locally uniformly on $\mathbb C$.
\end{definition}

\begin{lstlisting}
def StripClass (σ : ℝ) (f : ℂ → ℂ) : Prop :=
  ∃ P : ℕ → ℝ[X], (∀ n, P n ≠ 0) ∧ (∀ n (w : ℂ), aeval w (P n) = 0 → w.im ^ 2 ≤ σ) ∧
    TendstoLocallyUniformly (fun n z => aeval z (P n)) f atTop
\end{lstlisting}

Write $S_h f(z)=f(z+ih)+f(z-ih)$ for the imaginary shift operator.

\begin{proposition}[\texttt{StripClass.shift}, \texttt{StripClass.shiftIter}]\label{prop:shift}
Let $\sigma\ge0$, $h>0$.  If $f$ lies in the strip class at $\sigma$ then $S_hf$ lies in the strip
class at $\max(\sigma-h^{2},0)$; consequently $S_h^{N}f$ lies in the strip class at
$\max(\sigma-Nh^{2},0)$.
\end{proposition}

\begin{proof}
The polynomial witnesses transport: $S_h$ applied to a real polynomial is again a real polynomial,
because the complex polynomial $p(X+ih)+p(X-ih)$ has real coefficients (the odd-order Hasse
derivatives cancel; \texttt{shiftPolyC\_coeff\_real}), and its root strip is controlled by
Proposition~\ref{prop:polystep}.  Local uniform convergence is preserved because precomposition with
$z\mapsto z\pm ih$ preserves it and sums of locally uniformly convergent sequences converge locally
uniformly.  Nonvanishing of the shifted witness is \texttt{shiftPoly\_ne\_zero}.  The iterate is an
induction on $N$, whose only content is the arithmetic
$\max(\max(\sigma-Nh^{2},0)-h^{2},0)=\max(\sigma-(N{+}1)h^{2},0)$.
\end{proof}

\section{Transversality: the shift crosses, it never grazes}\label{sec:transversal}

The narrowing step of \S\ref{sec:class} controls \emph{where} the zeros of $S_hf$ are.  A separate and
equally elementary argument controls \emph{how} they are met: at the polynomial level the shift's real
zeros are always simple.  This is not needed for Theorem~\ref{thm:main}, and we record it because it
is the exact point at which multiplicity enters the construction, and because it is sharp
(Remark~\ref{rem:sharp}).

The observation that starts it: $p$ has real coefficients, so for real $x$ and real $h$,
$p(x-ih)=\overline{p(x+ih)}$, whence on the axis
\begin{equation}\label{eq:axisshift}
  (S_hp)(x)=2\Reop p(x+ih),\qquad (S_hp)'(x)=2\Reop p'(x+ih).
\end{equation}
So a \emph{double} zero of $S_hp$ at $x$ says exactly that both $p(w)$ and $p'(w)$ are purely
imaginary at $w=x+ih$ --- and then $\overline{p(w)}\,p'(w)$ is real.  It never is:

\begin{lemma}[\texttt{im\_conj\_mul\_deriv\_neg}]\label{lem:herglotz}
Let $p\in\mathbb R[X]$ have positive degree and only real zeros, and let $\Imop w>0$.  Then
\[
  \Imop\bigl(\overline{p(w)}\,p'(w)\bigr)<0 .
\]
\end{lemma}

The reason is the classical one --- $\overline{p}p'/|p|^{2}=\sum_k 1/(w-r_k)$ and each term has
imaginary part $-\Imop w/|w-r_k|^{2}<0$ --- but the formalization does not construct that sum.  It
peels one real root at a time: writing $p=(X-r)q$ with $r$ real and $q$ again real-rooted,
\begin{equation}\label{eq:peel}
  \Imop\bigl(\overline{p(w)}p'(w)\bigr)
  \;=\;-\,|q(w)|^{2}\,\Imop w\;+\;|w-r|^{2}\,\Imop\bigl(\overline{q(w)}q'(w)\bigr),
\end{equation}
in which the first term is strictly negative --- a real-rooted $q$ cannot vanish off the axis, so
$|q(w)|^{2}>0$ --- and the second is $\le0$ by induction on the degree.  The strict sign is produced by
the peel itself, which is why no sum and no partial-fraction decomposition is required.

\begin{theorem}[\texttt{shiftPoly\_zero\_simple}]\label{thm:simple}
Let $p\in\mathbb R[X]$ have positive degree and only real zeros, and let $h>0$.  Then every real zero
of $S_hp$ is simple.
\end{theorem}

\begin{proof}
At a zero, \eqref{eq:axisshift} gives $\Reop p(w)=0$; if it were not simple we would also have
$\Reop p'(w)=0$, making $\overline{p(w)}p'(w)$ real and contradicting Lemma~\ref{lem:herglotz}.
\end{proof}

The statement is carried to the compiled shift polynomial by
$\texttt{shiftPoly\_derivative}:\ \partial(S_hp)=S_h(\partial p)$ --- the shift commutes with
differentiation --- so Theorem~\ref{thm:simple} is a statement about
\texttt{derivative (shiftPoly h p)} and not merely about the shifted evaluations.

\begin{remark}[why this does \emph{not} descend to $\Xi_\lambda$]\label{rem:nodescend}
Theorem~\ref{thm:simple} is a statement about each approximating polynomial, and simplicity is not a
closed condition: Hurwitz (\S\ref{sec:hurwitz}) carries real-rootedness to a locally uniform limit but
not simplicity, since a sequence of simple zeros may coalesce.  One might hope to recover it from a
uniform lower bound on the zero gaps of the approximants, but no such bound is available at this
level: the approximants have degree $\to\infty$, and the heat flow applied to $z^{n}$ produces the
Hermite polynomial, whose $n$ roots spread over a range $\sim\sqrt n$ so that the minimum gap
degrades.  Any proof of simplicity for $\Xi_\lambda$ must therefore draw on the limit's own zero
density rather than on the polynomial step.  We do not prove it.
\end{remark}

\begin{remark}[the threshold is sharp for simplicity]\label{rem:sharp}
Simplicity cannot hold at $\lambda=\tfrac18$ itself.  The isolated pair $p(z)=z^{2}+\tfrac14$ flows to
$p_t(z)=z^{2}+\tfrac14-2t$, which at $t=\tfrac18$ is exactly $z^{2}$ --- a genuine double zero --- and
for $t>\tfrac18$ has the two distinct real roots $\pm\sqrt{2t-\tfrac14}$.  So the correct statement to
aim at is simplicity on the \emph{open} range $\lambda>\tfrac18$, matching the classical assertion that
the zeros of $\Xi_\lambda$ are real for $\lambda\ge\tfrac18$ and simple for $\lambda>\tfrac18$.
Numerically the separation is immediate rather than gradual: applying the flow to real-rooted
polynomials with multiple roots ($z^{2}$, $z^{3}$, $z^{4}$, $(z^2)(z-1)^2$, $(z+1)^2z(z-1)^2$) gives
all-real all-simple roots already at $t=10^{-8}$, with minimum gap growing like $2\sqrt{2t}$.
\end{remark}

\section{From shifts to the heat flow}\label{sec:cosh}

On the Fourier side the shift is a multiplication operator:
\[
  S_h\!\left(\int K(u)e^{izu}du\right)
  =\int K(u)\bigl(e^{-hu}+e^{hu}\bigr)e^{izu}du
  =\int 2\cosh(hu)\,K(u)\,e^{izu}\,du .
\]
Iterating, $S_h^{N}\Xi(z)=\int\bigl(2\cosh(hu)\bigr)^{N}K(u)e^{izu}du$; this is
\texttt{shiftOp\_iterate\_Ktrans}.  Now choose
\[
  h_n=\sqrt{\frac{2\lambda}{n+1}},\qquad N=n+1,
  \qquad\text{so that}\qquad Nh_n^{2}=2\lambda\ \text{ for every } n,
\]
and normalize by $2^{-(n+1)}$.  Since
$\cosh(h_nu)^{n+1}\ \text{is squeezed between}\ \bigl(1+\tfrac{\lambda u^{2}}{n+1}\bigr)^{n+1}$ and
$e^{\lambda u^{2}}$ --- the lower bound from $1+r^{2}/2\le\cosh r$, the upper from
$\cosh r\le e^{|r|}$ --- dominated convergence gives:

\begin{proposition}[\texttt{tendstoLocallyUniformly\_coshApprox}]\label{prop:cosh}
For $0\le\lambda\le1$, the normalized stages
$\;z\mapsto\int\cosh(h_nu)^{n+1}K(u)e^{izu}du\;$ converge to $\Xi_\lambda$ locally uniformly on
$\mathbb C$.
\end{proposition}

The dominating function is uniform on each horizontal strip $|\Imop z|\le A$, which is what
``locally uniformly on $\mathbb C$'' requires; the compact-set formulation is obtained from
\texttt{tendstoLocallyUniformly\_iff\_forall\_isCompact}.

Two nonvanishing facts are needed and are proved separately: no weighted transform
$\int w(u)K(u)e^{izu}du$ with $w$ continuous, positive and of at most exponential growth is
identically zero (\texttt{coshStage\_ne\_zero}, \texttt{flowStage\_ne\_zero}).  These follow from
Fourier inversion together with $\xi(\tfrac12)\ne0$.  They are exactly the hypothesis ``$f$ is not
identically zero'' of the Hurwitz step, and omitting them would make the argument vacuous rather than
wrong, which is why we flag them.

\section{\texorpdfstring{$\Xi$ lies in the strip class at $\tfrac14$}{Xi lies in the strip class at 1/4}}\label{sec:hadamard}

This is the arithmetic input, and the only place where anything about $\zeta$ beyond \eqref{eq:strip}
is used.

\begin{proposition}[\texttt{xiSection\_stripClass}]\label{prop:strip}
$\Xi$ lies in the strip class at $\sigma=\tfrac14$.
\end{proposition}

\begin{proof}[Construction of the witnesses]
Start from the Hadamard product, compiled as \texttt{riemannXi\_hadamard\_factorization}:
\[
  \xi(s)=e^{As+B}\prod_{\rho}\Bigl(1-\frac s\rho\Bigr)^{m_\rho}e^{m_\rho s/\rho},
\]
the product over nontrivial zeros with multiplicity.  Restrict to the cofinal finite sets
$G_N=\{\rho:|\Imop\rho|\le N\}$, which are closed under $\rho\mapsto1-\rho$ and $\rho\mapsto\bar\rho$
by the functional equation and the reflection $\overline{\xi(\bar s)}=\xi(s)$.  Substituting
$s=\tfrac12+iz$ and writing $w_\rho=-i(\rho-\tfrac12)$,
\[
  e^{As+B}\prod_{\rho\in G_N}\Bigl(\bigl(1-\tfrac s\rho\bigr)e^{s/\rho}\Bigr)^{m_\rho}
  \;=\;\kappa_N\,e^{iz(A+c_N)}\,Q_N(z),
  \qquad
  Q_N(z)=\prod_{\rho\in G_N}\Bigl(1-\frac z{w_\rho}\Bigr)^{m_\rho},
\]
with $c_N=\sum_{G_N}m_\rho/\rho$.  Because $G_N$ is closed under the two involutions, the zero
multiset of $Q_N$ is closed under conjugation, so $Q_N\in\mathbb R[X]$; and $w_\rho$ satisfies
$|\Imop w_\rho|=|\Reop\rho-\tfrac12|<\tfrac12$ by \eqref{eq:strip}, so every root of $Q_N$ lies in
$(\Imop w)^{2}\le\tfrac14$.  The paired sums converge, $c_N\to-A$
(\texttt{reflection\_const\_eq\_zero}), so the exponential factor tends to $1$ and
$\kappa_N\to\Xi(0)\ne0$.  Hence $\Xi(0)\,Q_N\to\Xi$ locally uniformly, and after rescaling these are
the required witnesses.
\end{proof}

The cancellation $c_N\to-A$ is the one delicate point: individually the sums $\sum1/\rho$ do not
converge absolutely, and it is the pairing $\rho\leftrightarrow1-\rho$ inside each $G_N$ that makes
$c_N$ converge.  This is why the exhaustion must be by symmetric sets and not by, say, the first $N$
zeros in some order.

\section{Real zeros survive locally uniform limits}\label{sec:hurwitz}

\begin{proposition}[\texttt{Hurwitz.im\_eq\_zero\_of\_tendstoLocallyUniformly}]\label{prop:hurwitz}
Let $F_n$ be entire with only real zeros, $F_n\to f$ locally uniformly on $\mathbb C$, and $f$ not
identically zero.  Then $f$ has only real zeros.
\end{proposition}

The formalization proves this through a nonvanishing form of Hurwitz's theorem on an open preconnected
set, and it is worth recording the route because it avoids winding numbers entirely.

\begin{lemma}[\texttt{Hurwitz.ne\_zero\_of\_tendstoLocallyUniformlyOn}]\label{lem:hurwitz}
Let $U\subseteq\mathbb C$ be open and preconnected, $F_n$ holomorphic on $U$ with $F_n$ nowhere zero on
$U$, $F_n\to f$ locally uniformly on $U$, and $f\not\equiv0$ on $U$.  Then $f$ is nowhere zero on $U$.
\end{lemma}

\begin{proof}
Suppose $f(z_0)=0$, $z_0\in U$.  Since $f\not\equiv0$ on the preconnected $U$, the identity theorem
gives a punctured neighbourhood of $z_0$ on which $f\ne0$; choose $r>0$ with
$\overline{B}(z_0,r)\subseteq U$ and $f\ne0$ on the sphere $|z-z_0|=r$, and let
$m=\min_{|z-z_0|=r}\|f(z)\|>0$ (attained, the sphere being compact).  Local uniform convergence is
uniform on $\overline{B}(z_0,r)$, so for some $n$ we have $\|F_n-f\|<m/2$ there, whence
\[
  \|F_n(z)\|>m/2 \ \text{ on the sphere},\qquad \|F_n(z_0)\|<m/2 \ \text{ at the centre}.
\]
But $F_n$ is nowhere zero on the disc, so $1/F_n$ is holomorphic there, and the maximum modulus
principle bounds $\|1/F_n(z_0)\|$ by its maximum on the boundary, i.e.\ $\le2/m$ --- that is,
$\|F_n(z_0)\|\ge m/2$, contradicting the second display.
\end{proof}

\begin{proof}[Proof of Proposition~\ref{prop:hurwitz}]
Apply Lemma~\ref{lem:hurwitz} to the two open half-planes $U_\pm=\{\pm\Imop w>0\}$.  Each is convex,
hence preconnected; each $F_n$ is nowhere zero on $U_\pm$ precisely because its zeros are all real; and
$f\not\equiv0$ on $U_\pm$ because $f\not\equiv0$ on $\mathbb C$ and $f$ is entire, so vanishing on an
open subset would force vanishing identically.  Hence $f$ has no zero off the real axis.
\end{proof}

Specialised to the strip class this gives \texttt{StripClass.im\_eq\_zero\_of\_zero}: a member of the
strip class at $\sigma=0$ which is not identically zero has only real zeros.

\section{Assembly}\label{sec:assembly}

\begin{proof}[Proof of Theorem~\ref{thm:main}]
Fix $\lambda$ with $\tfrac18\le\lambda\le1$ and let $n\in\mathbb N$.  By
Proposition~\ref{prop:strip}, $\Xi$ lies in the strip class at $\tfrac14$.  Apply
Proposition~\ref{prop:shift} with $h=h_n$ and $N=n+1$: the iterate $S_{h_n}^{n+1}\Xi$ lies in the
strip class at
\[
  \max\!\left(\tfrac14-(n{+}1)h_n^{2},\,0\right)=\max\!\left(\tfrac14-2\lambda,\,0\right)=0,
  \qquad\text{since }\lambda\ge\tfrac18 .
\]
By \texttt{shiftOp\_iterate\_Ktrans} that iterate is the cosh stage, and by
Proposition~\ref{prop:hurwitz} (applied within the strip class, using
\texttt{coshStage\_ne\_zero}) every zero of the $n$-th normalized stage is real.  Now let
$n\to\infty$: by Proposition~\ref{prop:cosh} the stages converge to $\Xi_\lambda$ locally uniformly,
and $\Xi_\lambda\not\equiv0$ by \texttt{flowStage\_ne\_zero}, so Proposition~\ref{prop:hurwitz} applies
once more and every zero of $\Xi_\lambda$ is real.
\end{proof}

Hurwitz is used twice, at two different levels --- once to pass from polynomial witnesses to each
fixed stage, once to pass from the stages to the flow.  The threshold $\tfrac18$ enters at exactly one
place: the requirement $2\lambda\ge\sigma=\tfrac14$.

\begin{corollary}[\texttt{flowCascade\_realRooted\_of\_le\_seven}]
In the cascade $\lambda_n=1/(n+1)$, the stages $n\le7$ have only real zeros, at every height.
\end{corollary}

\section{The formalization}\label{sec:lean}

The development is Lean~4 against Mathlib \cite{Mathlib} at the revision pinned in \texttt{lake-manifest.json}
(\texttt{v4.33.0-rc2}).

\begin{center}
\begin{tabular}{lrrl}
\toprule
module & lines & thms & content \\
\midrule
\texttt{XiKernelTheta}                & 123 & 16 & $\Psi$ holomorphic on a strip \\
\texttt{XiKernelEven}                 & 122 & 11 & $K(-u)=K(u)$ from the theta reflection \\
\texttt{XiKernelTransform}            & 269 & 22 & $\int K(u)e^{izu}du=\Xi(z)$ \\
\texttt{XiHadamardFactorization}      & 484 & 13 & the Hadamard product for $\xi$ \\
\texttt{XiProductMultPartialFraction} & 277 &  6 & multiplicity bookkeeping, $\sum m_\rho/\rho$ \\
\texttt{XiStripApproximants}          & 542 & 43 & Prop.~\ref{prop:strip}: $\Xi$ in the class at $\tfrac14$ \\
\texttt{DeBruijnStripStep}            & 230 & 13 & Lem.~\ref{lem:core}, Prop.~\ref{prop:polystep} \\
\texttt{HurwitzLimit}                 & 139 &  2 & Prop.~\ref{prop:hurwitz} \\
\texttt{StripShiftClass}              & 166 & 14 & \texttt{StripClass}, Prop.~\ref{prop:shift} \\
\texttt{FiberShiftSimple}            & 226 &  7 & Lem.~\ref{lem:herglotz}, Thm.~\ref{thm:simple} \\
\texttt{XiFlowCascade}                & 366 & 30 & $\Xi_\lambda$, the cascade $\lambda_n$ \\
\texttt{FlowCoshStages}               & 354 & 26 & Prop.~\ref{prop:cosh}, nonvanishing \\
\texttt{FlowFrontier}                 &  69 &  5 & \texttt{FlowFrontier}, monotonicity \\
\texttt{DeBruijnFlowFrontier}         &  95 &  8 & Theorem~\ref{thm:main} \\
\midrule
total                                 & 3462 & 216 & \\
\bottomrule
\end{tabular}
\end{center}

\paragraph{Axiom audit.}  Every named theorem in the chain, checked with \verb|#print axioms|,
depends on exactly $\{\texttt{propext},\ \texttt{Classical.choice},\ \texttt{Quot.sound}\}$ --- the
three axioms of Lean's standard classical foundation --- and nothing else.  There is no \texttt{sorry}
and no bespoke \texttt{axiom} in the development or anywhere in the repository containing it: a full
\verb|lake build| over $10451$ jobs reports zero errors and zero declarations using \texttt{sorry}.

\paragraph{What Mathlib supplied and what it did not.}  Mathlib supplied the analytic infrastructure:
holomorphy, locally uniform convergence, the maximum modulus principle, the identity theorem
(\texttt{eqOn\_zero\_of\_preconnected\_of\_eventuallyEq\_zero}), dominated convergence, Fourier
integrals, the Jacobi theta function with its functional equation
(\texttt{evenKernel\_functional\_equation}), polynomial algebra over $\mathbb R$ and $\mathbb C$
including \texttt{aroots} with
multiplicity, and \texttt{completedRiemannZeta} with its functional equation.  Mathlib did \emph{not}
supply, and this development therefore builds from scratch: the Hadamard product for $\xi$; any notion
of a real-rooted or Laguerre--P\'olya class; preservation of real-rootedness under locally uniform
limits; and any de Bruijn--Newman material whatsoever.  A search of the pinned revision for
\texttt{Laguerre}, \texttt{P\'olya}, \texttt{realRooted} and \texttt{deBruijn} returns nothing relevant;
\texttt{Mathlib/Analysis/Complex/Hadamard.lean} is the three-lines theorem, not the factorization.

\section{The remaining interval}\label{sec:remaining}

Let \texttt{FlowFrontier} $\lambda_0$ abbreviate ``every $\Xi_\lambda$ with
$\lambda_0\le\lambda\le1$ has only real zeros''.  The following is also compiled:

\begin{lstlisting}
theorem rh_of_flowFrontier (h : ∀ lam₀ : ℝ, 0 < lam₀ → FlowFrontier lam₀) : RiemannHypothesis
\end{lstlisting}

\noindent Its hypothesis quantifies over \emph{every} positive $\lambda_0$; the proof instantiates at
$\lambda_n=1/(n+1)\to0$.  Theorem~\ref{thm:main} supplies $\lambda_0\ge\tfrac18$.  The gap is therefore
the bounded interval
\[
  \lambda\in\left(0,\tfrac18\right),
\]
which under \eqref{eq:norm} is the statement $\Lambda\le0$, and by Newman's conjecture \cite{Newman1976}, settled by
Rodgers--Tao ($\Lambda\ge0$) \cite{RodgersTao2020}, this is equivalent to the Riemann Hypothesis.  Nothing in this paper
bears on it.

It is worth recording precisely \emph{where} the constant $\tfrac18$ comes from, because it is not a
feature of $\Xi$.  Proposition~\ref{prop:polystep} is exactly tight for a single conjugate pair and
nothing else: for $p(z)=z^{2}+y_0^{2}$ one has $p''=2$, so under the flow
$p_t(z)=z^{2}+y_0^{2}-2t$, whose roots reach the real axis exactly at $t=y_0^{2}/2$; with
$y_0=\tfrac12$ this is $\tfrac18$.  The input $\sigma=\tfrac14$ of Proposition~\ref{prop:strip} is the
bare critical strip \eqref{eq:strip}: no zero density, no zero count, no arithmetic beyond
$0<\beta<1$ enters Theorem~\ref{thm:main} at all.

\begin{remark}[a numerical experiment, not a theorem]\label{rem:density}
The flow acts exactly on polynomials, $p_t=\sum_j(-t)^{j}p^{(2j)}/j!$, and for a conjugate pair at
$x\pm iy$ in a field of further zeros $x_k$ the imaginary part obeys
\[
  \partial_t y \;=\; -\frac1y \;-\; 2y\sum_k\frac1{(x-x_k)^{2}+y^{2}} ,
\]
in which \emph{both} terms are negative when the $x_k$ are real: neighbours accelerate the descent, so
$t^{*}\le y_0^{2}/2$ should be slack.  We measured this.  With $y_0=\tfrac12$ and $m$ real zeros at
spacing $s=1/\rho$, the first time $t^{*}$ at which all roots are real satisfies (control: $m=0$
reproduces $t^{*}=\tfrac18$ to $8\times10^{-23}$)
\[
  \begin{array}{c|ccc}
    \rho & 1 & 2 & 4\\\hline
    t^{*}/(1/8)\ \ (m=40) & 0.471 & 0.279 & 0.156
  \end{array}
\]
and at $\rho=4$ the value $0.0195$ agrees to $2\%$ with the linear law $y_0/(2\pi\rho)=0.0199$ that
replaces the quadratic $y_0^{2}/2$ once the density is felt.  Since the zeros of $\Xi$ have density
$\rho=\log(T/2\pi)/2\pi$ at height $T$, the corresponding heuristic threshold is
$y_0/\log(T/2\pi)\to0$.  Replacing all neighbours by conjugate pairs at the maximal height
$y_1=y_0=\tfrac12$ --- the configuration best able to push the pair \emph{away} from the axis, and the
relevant one since the neighbours of a hypothetical off-line zero are not known to be real --- the
ratios become $0.680$, $0.479$, $0.329$: still decreasing in $\rho$.  De Bruijn's bound held as a
control throughout ($t^{*}\le\tfrac18$ in every configuration tested).

We draw no conclusion about $\Xi$ from this.  What it does identify is the point at which a
sharpening would have to enter: a \emph{local, density-aware} refinement of
Proposition~\ref{prop:shift}, in which $m$ roots in a window of width $w$ about $x$ buy a narrowing
near $x$ strictly better than $Nh^{2}$, by an amount increasing in $m/w$.  That is a statement about
polynomials.
\end{remark}

\section{What is not claimed}

The theorem proved here is Theorem~\ref{thm:main} and nothing more.  In particular:

\begin{itemize}
\item The Riemann Hypothesis is not proved, and no step of this development is a step toward it.  The
  distance between Theorem~\ref{thm:main} and RH is the interval $\lambda\in(0,\tfrac18)$, and closing
  that interval \emph{is} RH; it is not a gap in formalization or in plumbing.
\item Theorem~\ref{thm:main} is de Bruijn's 1950 bound \cite{deBruijn1950} and is weaker than the
  published state of the art.  Ki--Kim--Lee sharpened it to the strict inequality $\Lambda<\tfrac12$
  \cite{KKL2009}, and Polymath~15 gave $\Lambda\le0.22$ \cite{Polymath15}, i.e.\ $\lambda\le0.055$,
  roughly a factor $2.3$ better; neither is formalized here, and formalizing Polymath~15 would shrink
  the open interval to $(0,0.055)$ without closing it.
\item The novelty claimed is formalization, not mathematics.  The mathematical content of
  \S\ref{sec:core}--\S\ref{sec:assembly} is de Bruijn's, and where the argument is arranged differently
  from his --- the strip class as an explicit predicate, the $\cosh$-power staging, the symmetric
  exhaustion $G_N$ --- that is a convenience of formalization, not a new idea.
\item Simplicity of the zeros of $\Xi_\lambda$ is \emph{not} proved.  Theorem~\ref{thm:simple} is a
  statement about the approximating polynomials only; Remark~\ref{rem:nodescend} explains why it does
  not descend through the limit, and Remark~\ref{rem:sharp} why any such statement must be confined to
  the open range $\lambda>\tfrac18$.  A fortiori nothing here bears on the simplicity of the zeros of
  $\zeta$, which would be a statement at $\lambda=0$.
\item Remarks~\ref{rem:sharp} and \ref{rem:density} report numerical experiments on polynomials, with
  pre-registered decision rules and positive controls.  They are not theorems and are not used
  anywhere above.
\end{itemize}

\section*{Availability}

The development builds with \texttt{lake build} against the pinned Mathlib.  The two results stated
above are
\begin{center}
\begin{tabular}{ll}
Theorem~\ref{thm:main} & \texttt{CriticalLinePhasor.XiKernel.flowFrontier\_one\_eighth} \\
  & in \texttt{RequestProject/DeBruijnFlowFrontier.lean} \\[2pt]
Theorem~\ref{thm:simple} & \texttt{CriticalLinePhasor.DeBruijn.shiftPoly\_zero\_simple} \\
  & in \texttt{RequestProject/FiberShiftSimple.lean}
\end{tabular}
\end{center}

\begin{thebibliography}{9}

\bibitem{deBruijn1950}
N.~G.~de Bruijn,
\emph{The roots of trigonometric integrals},
Duke Math.\ J.\ \textbf{17} (1950), 197--226.

\bibitem{Newman1976}
C.~M.~Newman,
\emph{Fourier transforms with only real zeros},
Proc.\ Amer.\ Math.\ Soc.\ \textbf{61} (1976), 245--251.

\bibitem{KKL2009}
H.~Ki, Y.-O.~Kim, J.~Lee,
\emph{On the de Bruijn--Newman constant},
Adv.\ Math.\ \textbf{222} (2009), 281--306.

\bibitem{Polymath15}
D.~H.~J.~Polymath,
\emph{Effective approximation of heat flow evolution of the Riemann $\xi$ function, and a new upper
bound for the de Bruijn--Newman constant},
Res.\ Math.\ Sci.\ \textbf{6} (2019), no.~3, Paper No.~31.

\bibitem{RodgersTao2020}
B.~Rodgers, T.~Tao,
\emph{The de Bruijn--Newman constant is non-negative},
Forum Math.\ Pi \textbf{8} (2020), e6.

\bibitem{Titchmarsh}
E.~C.~Titchmarsh,
\emph{The Theory of the Riemann Zeta-Function}, 2nd ed., rev.\ D.~R.~Heath-Brown,
Oxford University Press, 1986.

\bibitem{Mathlib}
The mathlib Community,
\emph{The Lean mathematical library},
CPP 2020, 367--381.

\end{thebibliography}

\end{document}
